```latex
% ch2_algorithmes.tex

\chapter{Algorithmes}
\label{ch:algorithmes}

StreamKM++ est un algorithme de clustering conçu pour traiter des flux de données dont la taille peut rendre impossible la conservation de l'ensemble des observations en mémoire. Son principe consiste à construire progressivement un \emph{coreset}, c'est-à-dire un résumé pondéré et de taille limitée des données déjà observées. Le clustering final est ensuite effectué sur ce résumé plutôt que sur l'ensemble du flux \cite{ackermann2012}.

L'approche repose sur plusieurs composants complémentaires : une version pondérée de l'algorithme de Lloyd, l'initialisation k-means++, la construction d'un coreset à l'aide d'un arbre de partitionnement et une stratégie \emph{merge-and-reduce} permettant de maintenir ce résumé au fur et à mesure de l'arrivée des données. Ces composants sont présentés dans les sections suivantes, en précisant les choix effectués dans l'implémentation lorsqu'ils diffèrent de la formulation de référence.

% ---------------------------------------------------------------------------
\section{Algorithme de Lloyd pondéré}
% ---------------------------------------------------------------------------

L'algorithme de Lloyd constitue l'étape d'optimisation utilisée par k-means. À partir d'un ensemble de centres initiaux, il alterne deux opérations jusqu'à convergence. Chaque point est d'abord affecté au centre le plus proche, puis chaque centre est recalculé à partir des points qui lui sont associés.

Dans StreamKM++, cette opération est appliquée au coreset plutôt qu'aux observations originales. Les représentants du coreset possèdent un poids correspondant à la quantité de données qu'ils représentent. Ces poids doivent donc être pris en compte lors du recalcul des centres.

Soit un ensemble de points pondérés

\[
P = \{(p_i,w_i)\}_{i=1}^{n}
\]

et un ensemble de centres

\[
C = \{c_1,\ldots,c_k\}.
\]

Après l'affectation de chaque point à son centre le plus proche, le centre du cluster \(j\) est recalculé selon :

\[
c_j =
\frac{
    \sum_{p_i \in C_j} w_i p_i
}{
    \sum_{p_i \in C_j} w_i
}.
\]

La qualité du clustering est évaluée à l'aide de la somme pondérée des distances euclidiennes au carré :

\[
\mathrm{cost}(P,C)
=
\sum_{i=1}^{n}
w_i
\min_{c \in C}
\|p_i-c\|^2.
\]

La pondération permet à un représentant de poids \(w\) d'avoir le même effet sur le calcul des centres que \(w\) observations identiques de poids unitaire. Elle est donc indispensable lorsqu'un point du coreset représente plusieurs observations du jeu de données initial.

Dans l'implémentation réalisée, cette propriété est vérifiée par des tests unitaires. Lorsqu'un cluster devient vide au cours d'une itération, son centre précédent est conservé afin d'éviter une réinitialisation supplémentaire. L'algorithme s'arrête lorsque l'amélioration relative du coût devient inférieure à un seuil fixé ou lorsque le nombre maximal d'itérations est atteint.

% ---------------------------------------------------------------------------
\section{Initialisation k-means++}
% ---------------------------------------------------------------------------

L'algorithme de Lloyd est sensible au choix des centres initiaux. Une sélection entièrement aléatoire peut conduire à des centres concentrés dans une même région des données et ainsi dégrader la qualité du résultat obtenu.

L'initialisation k-means++ proposée par Arthur et Vassilvitskii \cite{arthur2007} cherche à mieux répartir les centres initiaux. Après la sélection d'un premier centre, chaque nouveau centre est choisi avec une probabilité proportionnelle au carré de sa distance au centre déjà sélectionné le plus proche.

Pour un point \(p\), on définit :

\[
D(p)^2 =
\min_{c \in C}
\|p-c\|^2.
\]

Dans le cas pondéré utilisé avec le coreset, la probabilité de sélectionner un point \(p_i\) comme nouveau centre est proportionnelle à :

\[
w_i D(p_i)^2.
\]

Ainsi, les représentants ayant un poids important et situés loin des centres déjà sélectionnés ont une probabilité plus élevée d'être choisis.

\paragraph{Choix du premier centre.}

Dans la formulation classique de k-means++, le premier centre est sélectionné uniformément parmi les observations \cite{arthur2007}. Dans l'implémentation réalisée, lorsque les données sont pondérées, ce premier centre est choisi avec une probabilité proportionnelle au poids du point. Les deux approches sont équivalentes lorsque tous les poids valent \(1\). Dans le cas d'un coreset, ce choix permet de prendre en compte dès l'initialisation la masse de données représentée par chaque point.

Afin de limiter la sensibilité du résultat à l'initialisation, le processus composé de k-means++ puis de l'algorithme de Lloyd est exécuté plusieurs fois. Le modèle présentant le coût final le plus faible est conservé. Dans les expériences réalisées, cinq exécutions sont utilisées sauf indication contraire.

% ---------------------------------------------------------------------------
\section{Construction du coreset}
% ---------------------------------------------------------------------------

L'objectif principal de StreamKM++ est de réduire la quantité de données conservée tout en maintenant suffisamment d'information pour obtenir un clustering représentatif. Cette réduction repose sur la construction d'un coreset \cite{ackermann2012}.

Un coreset est un ensemble pondéré de taille réduite représentant les principales caractéristiques des données originales. Chaque point du coreset est associé à un poids correspondant à la masse des observations qu'il résume.

Dans StreamKM++, ce résumé est construit à l'aide d'un \emph{coreset tree}. Les données sont organisées progressivement dans un arbre binaire dont chaque feuille correspond à un sous-ensemble des observations et possède un représentant.

Initialement, une seule feuille contient l'ensemble des points à résumer. À chaque étape, une feuille est sélectionnée pour être subdivisée. La probabilité de sélection dépend de son coût : une région dont les données sont mal représentées possède un coût plus important et est donc davantage susceptible d'être raffinée.

Une fois la feuille sélectionnée, un nouveau représentant est choisi à partir des distances au représentant existant. Les observations sont ensuite réparties entre les deux représentants selon leur proximité. Ce processus est répété jusqu'à atteindre la taille cible du coreset ou jusqu'à ce qu'aucune subdivision supplémentaire ne soit possible.

\paragraph{Coût et masse.}

Deux grandeurs doivent être distinguées lors de la construction de l'arbre. Le coût d'un nœud est calculé à partir des distances entre les observations qu'il contient et son représentant. Il intervient dans le choix des régions à subdiviser.

La masse d'un nœud correspond quant à elle à la somme des poids des observations qu'il représente. Elle devient le poids du représentant lorsque celui-ci est intégré au coreset.

Dans l'implémentation réalisée, le coût et la masse sont donc représentés séparément. Le premier intervient dans la construction de l'arbre, tandis que la seconde est utilisée pour pondérer les représentants produits.

\paragraph{Conservation de la masse.}

La construction du coreset doit préserver la masse totale des données. Si l'ensemble d'entrée possède une masse

\[
W =
\sum_i w_i,
\]

alors le coreset \(S\) doit vérifier :

\[
\sum_{q \in S} w_q = W.
\]

La réduction du nombre de représentants ne modifie donc pas la quantité totale de données représentée. Cette propriété constitue l'un des invariants vérifiés par les tests de l'implémentation.

\paragraph{Cas dégénérés.}

La valeur \(m\) définit une taille cible maximale pour le coreset. Dans le cas général, la construction cherche à produire \(m\) représentants. Certaines configurations peuvent toutefois empêcher une nouvelle subdivision, notamment lorsque plusieurs observations sont identiques et que le coût d'une feuille devient nul.

Dans cette situation, le coreset peut contenir moins de \(m\) représentants sans que cela constitue une erreur. Sa taille reste néanmoins bornée par \(m\) et la masse totale des données demeure conservée.

% ---------------------------------------------------------------------------
\section{Merge-and-reduce}
\label{sec:algo-merge-reduce}
% ---------------------------------------------------------------------------

Le \emph{coreset tree} permet de résumer un ensemble fini de points. Pour traiter progressivement un flux dont la taille totale n'est pas nécessairement connue à l'avance, StreamKM++ utilise une stratégie \emph{merge-and-reduce} \cite{ackermann2012}.

Le principe consiste à maintenir plusieurs niveaux de stockage, appelés buckets :

\[
B_0, B_1, B_2, \ldots
\]

Le premier niveau reçoit les observations du flux. Lorsque la quantité de données à traiter atteint le seuil défini par la taille du coreset, les données sont résumées puis propagées vers les niveaux supérieurs.

Si le niveau suivant est libre, le résumé y est stocké. Lorsqu'il contient déjà un coreset, les deux résumés sont réunis puis réduits à nouveau afin de produire un nouveau coreset de taille bornée. Cette propagation se poursuit tant que le niveau suivant est occupé.

Le fonctionnement est comparable à la propagation d'une retenue dans une représentation binaire. Les résumés sont progressivement fusionnés afin de représenter des portions de plus en plus importantes du flux sans conserver l'ensemble des observations originales.

\paragraph{Invariants de la structure.}

Les résumés stockés dans les différents niveaux respectent deux propriétés principales :

\begin{itemize}
    \item leur taille reste bornée par \(m\) ;
    \item la masse totale représentée est conservée lors des différentes fusions.
\end{itemize}

La taille d'un résumé peut être inférieure à \(m\) dans les cas dégénérés décrits précédemment. L'invariant porte donc sur une taille maximale plutôt que sur l'obtention systématique d'exactement \(m\) représentants.

\paragraph{Coreset courant.}

À tout instant, l'ensemble des buckets non vides constitue un résumé pondéré des observations déjà traitées. Le nombre de niveaux augmente avec la quantité de données reçue, ce qui permet de contrôler la consommation mémoire sans nécessiter la connaissance préalable de la taille totale du flux.

Lorsqu'un clustering est demandé, les représentants contenus dans les différents buckets sont regroupés afin de constituer le coreset courant. L'initialisation k-means++ et l'algorithme de Lloyd pondéré sont ensuite appliqués à ce résumé.

\paragraph{Fusion de résumés indépendants.}

La composition des coresets permet également de combiner des résumés construits indépendamment sur des sous-ensembles distincts des données. Les représentants issus de plusieurs résumés peuvent être réunis puis réduits à nouveau afin d'obtenir un coreset représentant leur union.

Cette propriété est particulièrement adaptée à l'exécution distribuée avec Spark. Chaque partition peut construire indépendamment son propre résumé à l'aide d'un \texttt{BucketManager}. Les coresets obtenus sont ensuite fusionnés progressivement afin de produire un résumé global.

La logique de StreamKM++ reste ainsi concentrée dans les packages \texttt{streamkm.core} et \texttt{streamkm.coreset}. La couche \texttt{streamkm.spark} prend en charge la distribution des traitements sans modifier le fonctionnement interne de la construction du coreset.
```
